% ============================================================================
% Chapter 22 --- Wave 4: Policy plane
% ----------------------------------------------------------------------------
% Purpose : install PDP/PEP, parameterised by identity attributes (Wave 1)
%           and observed by the audit pipeline (Wave 2). The most ambitious
%           novelty of the architecture for most universities (D7.5).
% Page target: ~5 pages.
% Dependencies: Wave 0 + Wave 1 + Wave 2 (D17.3); reference layer Ch. 7;
%               IdP attribute mapping per D7.7.
% Mirrors wave template (D18.2).
% ----------------------------------------------------------------------------
% Decisions registered in _coherence-EN.md:
%   D22.1 shadow mode is the mandatory introduction (D7.1)
%   D22.2 per-policy enforcement decision is the unit of progress
%   D22.3 brownfield single-suite-centric: conditional-access logic partially reproduced
%   D22.4 exit criteria reference Policy conformance subset (D16.5)
% ============================================================================

\chapter{Wave~4 --- Policy plane}
\label{ch:wave4-policy}

\begin{caveatbox}[Dependencies]
\noindent Wave~4 applies the wave chapter template established in
\trchap{ch:wave0-foundations}{}. Its required pre-requisites are
Waves 1 and 2: an identity layer that emits
attributes in the institutional schema, and an observability layer
that captures decisions. It instantiates the policy contract of
\trchap{ch:policy}{}, with the mapping of identity attributes to
the institutional schema specified in \trchap{ch:policy}{}. Closure is gated by
the four Policy tests of \trchap{ch:conformance}{}.
\end{caveatbox}

Wave~4 is, for most universities at the time of writing, the
most ambitious novelty of the playbook. Few institutions today
operate a dedicated policy decision point separate from their
applications and identity provider. Adopting one is a real
organisational change as well as a technical one: the institution
acquires a new artefact (the policy corpus), a new governance
process (policy authoring and review), and a new operational
capability (decision auditability). The wave's success depends as
much on the institutional change as on the technical deployment.

%-----------------------------------------------------------------------------
\section{Goal and dependencies}
\label{sec:w4-goal}

The goal of Wave~4 is to deploy a policy decision point conformant
to the contract of \trchap{ch:policy}{}, to migrate selected
authorisation decisions from the identity-provider-native rules and
from in-application logic into the policy plane, and to establish
the policy authoring and review process under which the corpus
will evolve.

\paragraph{Selected, not all.} Wave~4 does not require that every
authorisation decision in the institution be expressed in the
policy plane. The wave's success criterion is that the plane
exists, that it is conformant, that an institutionally agreed set
of decisions has been migrated, and that the process for migrating
further decisions is established. The transition of the rest of
the corpus is the work of subsequent waves and of routine
operations.

%-----------------------------------------------------------------------------
\section{Pre-requisites}
\label{sec:w4-prereq}

Beyond Waves 0, 1 and 2, three operational conditions strengthen
Wave~4:

\begin{itemize}[leftmargin=1.5em,nosep]
  \item the institutional attribute schema (D7.7) is stable; the
    identity layer of Wave~1 emits attributes against it;
  \item the existing authorisation logic --- whether expressed in
    the identity provider's rules, in application configuration,
    or in code --- is inventoried at the level of intent if not at
    the level of implementation;
  \item the security operations function (which will consume
    decision logs) is engaged with the canonical schema of
    \trchap{ch:observability}{}.
\end{itemize}

%-----------------------------------------------------------------------------
\section{Coexistence pattern}
\label{sec:w4-coexistence}

The principal engineering pattern of Wave~4 is \emph{shadow mode}:
the policy decision point evaluates every relevant request and
emits its decision into the audit trail, but the existing
enforcement points continue to apply the incumbent authorisation
logic. The institution thereby acquires the ability to compare
proposed decisions against actual behaviour before any decision is
moved to enforcing.

\paragraph{Per-policy transition.} Each authorisation rule is
migrated in three steps: it is expressed in the policy language;
the resulting decisions are observed in shadow mode for an
institutionally agreed period; if the discrepancy profile is
understood and the institution accepts the proposed behaviour, the
rule is switched to enforcing and the incumbent rule is retired.
The transition is per-policy, not per-application; an application
may have some of its authorisation in the policy plane and some
still in its incumbent location for a substantial period.

\paragraph{Shadow observation period.} An indicative period is two
to four weeks per policy for routine rules, longer for rules with
infrequent triggers (year-end administrative procedures,
research-data extraction patterns, incident-response invocations)
that may not be exercised within a shorter window. The observation
period is set per policy, not by a single institutional rule, and
is documented in the policy corpus alongside the policy itself.

%-----------------------------------------------------------------------------
\section{Trajectory: greenfield}
\label{sec:w4-greenfield}

For institutions in a greenfield position, Wave~4 establishes the
policy plane as the authorisation surface from the start: there is
no existing authorisation logic to migrate from. The wave is
correspondingly compressed but not trivial --- the policy corpus
must still be authored, the schema integration with identity
verified, and the auditing trail exercised. Indicative duration is
six to nine months.

%-----------------------------------------------------------------------------
\section{Trajectory: brownfield single-suite-centric}
\label{sec:w4-brownfield-ms}

For institutions whose authorisation logic resides primarily in
the incumbent identity suite's conditional-access (or equivalent)
expressions and in that suite's application-level configuration,
Wave~4 introduces the policy plane in front of, not in place of,
the existing authorisation arrangements.

\paragraph{Pattern.} A standalone policy decision point is
deployed alongside the incumbent identity suite. For each
conditional-access expression identified in Wave~1's inventory, an
equivalent policy is authored in the chosen policy language and
exercised in shadow mode. The incumbent suite's conditional-access
expression remains authoritative until the shadow comparison
demonstrates equivalence; at that point, the institution decides
whether to retire the vendor expression in favour of the
policy-plane rule or to keep it as a defence-in-depth element.

\paragraph{What the wave does not promise.} Wave~4 does not
promise that all conditional-access expressions will be migrated.
Some are unproductive to migrate (rules that genuinely belong with
identity, such as MFA enforcement); some are operationally costly
to migrate (rules that depend on vendor-specific signals not
exposed to the policy plane); some are simply not worth the
effort. The selection is institutional and is informed by the
sovereignty grid: rules that confer roadmap leverage to the
supplier are the priority candidates for migration.

\paragraph{Indicative duration.} Twelve to eighteen months,
dominated by the per-policy shadow observation and by the
institutional debate over which rules to migrate.

%-----------------------------------------------------------------------------
\section{Trajectory: brownfield hybrid}
\label{sec:w4-brownfield-hybrid}

For institutions whose authorisation logic is scattered across
standards-based federation attribute releases, open-source
identity-broker authorisation policies, in-application
configuration, and ad hoc code paths, Wave~4 is
principally the work of \emph{consolidation}.

\paragraph{Pattern.} The policy plane becomes the surface against
which scattered authorisation logic is gradually consolidated. The
three steps of the per-policy transition apply to each candidate
rule, regardless of where it currently resides. Rules that already
live in standards-aligned components (open-source identity-broker
authorisation policies expressed in standard ways) are migrated
more easily; rules that live in application code require more
engineering investment per migration.

\paragraph{What the wave changes.} The policy corpus becomes the
canonical home for institutional authorisation decisions. The
governance of the corpus --- review process, authoring rights,
deprecation discipline --- is established and exercised. The
decision log becomes a reliable forensic resource.

\paragraph{Indicative duration.} Twelve to eighteen months, with
substantial variation depending on the spread of the incumbent
authorisation logic and on the institution's appetite for code
refactoring.

%-----------------------------------------------------------------------------
\section{Reversibility criteria}
\label{sec:w4-reversibility}

Wave~4 is reversible at the per-policy level for as long as the
incumbent authorisation logic is preserved as defence-in-depth.
Beyond the point at which the incumbent rule is retired, reversal
is operationally possible but is treated as a re-authoring of the
incumbent rule rather than as a policy-plane reversal.

\begin{itemize}[leftmargin=1.5em,nosep]
  \item The shadow / enforcing toggle is per policy and is
    exercised through the policy bundle distribution mechanism;
    reversal is the deployment of a new bundle, not the
    redeployment of the engine.
  \item An emergency disable procedure --- a tested mechanism for
    switching the policy plane to permissive mode within minutes
    --- is part of the wave's deliverables.
  \item The full-wave rollback procedure (return to the
    pre-Wave-4 authorisation arrangements) is documented at wave
    open and exercised in tabletop form before any production
    enforcement migration.
\end{itemize}

%-----------------------------------------------------------------------------
\section{Exit criteria}
\label{sec:w4-exit}

Wave~4 closes when the four Policy tests of the conformance suite
of \trchap{ch:conformance}{} (\S\ref{sec:cf-suite}) pass, when the
agreed set of policies is enforcing, and when the governance
process for the corpus is operating.

\begin{itemize}[leftmargin=1.5em,nosep]
  \item PDP query API returns canonical decisions; recorded
    decisions replay identically; policy bundle export and
    re-import on alternative engine produces identical decisions
    on a representative input set; shadow-mode logging is at
    parity with enforcing-mode logging. (Policy conformance
    subset.)
  \item Migrated policy set: the institutionally agreed set of
    rules has been migrated and switched to enforcing; the
    remaining rules are either in a documented schedule or in the
    documented exception register.
  \item Governance: the policy corpus is under version control
    with a documented review process, and the first quarterly
    governance cycle has been completed.
\end{itemize}

\begin{realworldbox}[Policy plane adoption patterns]
\noindent The policy plane is more thinly documented in academia
than in industry. Industrial deployments of standalone policy
engines at scale (publicly described at organisations including
Netflix, Capital One and Pinterest) demonstrate that the
per-policy shadow-mode pattern
operates reliably under substantial load; their applicability to
academic settings is conditional on the institution's specific
authorisation profile. A cross-institutional academic network
would be one of the contexts in which a shared library of
academic-context policies might emerge.
\end{realworldbox}

%-----------------------------------------------------------------------------
\section{Regulatory anchoring}
\label{sec:w4-regulatory}

\noindent Wave~4 introduces the policy plane as the per-request
decision point. It contributes to the applicable data-protection
regime's data-protection-by-design obligation enforced at decision
time (e.g.\ \gls{gdpr}~Art.~25), and to the applicable
cybersecurity-baseline regime's risk-analysis-policy and
access-control-policy requirements (e.g.\ NIS2~Art.~21(2)(a) and
(i)). Under the applicable AI regulation it is the locus where the
prohibited-practices and transparency provisions (e.g.\ EU AI Act
Art.~5 and Art.~50) admit a
policy-plane enforcement: rules that gate AI-related
TACs on context labelling and on prohibited-use classification.
The wave operationalises 27001~A.5.15 (access control), A.5.18
(access rights) and A.8.3 (information access restriction) and
contributes to NIST CSF \textsf{PROTECT} (PR.AA). Detailed
mapping in \trchap{ch:regulatory-mapping}{} \S\ref{sec:rm-gdpr},
\S\ref{sec:rm-nis2}, \S\ref{sec:rm-ai-act},
\S\ref{sec:rm-iso}, \S\ref{sec:rm-nist-csf}.
